\documentclass[a4paper,12pt]{article}

\usepackage[T2A]{fontenc}
\usepackage[utf8]{inputenc}
\usepackage[russian]{babel}
\usepackage{amsmath,amssymb,mathtools}
\usepackage{helvet}
\renewcommand{\familydefault}{\sfdefault}
\usepackage{icomma}
\usepackage[a4paper,left=20mm,right=20mm,top=20mm,bottom=22mm,headheight=25pt]{geometry}
\usepackage{microtype}
\usepackage{tikz}
\usetikzlibrary{calc,arrows.meta,positioning,shapes.geometric}
\usepackage{tabularx,booktabs}
\usepackage{enumitem}
\usepackage{hyperref}
\usepackage{parskip}
\usepackage{fancyhdr}
\usepackage{setspace}
\usepackage{xcolor}

\definecolor{accent}{HTML}{008080}
\definecolor{darkaccent}{HTML}{004D4D}
\definecolor{textgray}{HTML}{333333}
\definecolor{softgray}{HTML}{6A6A6A}

\hypersetup{hidelinks}
\onehalfspacing
\setlength{\parindent}{0pt}
\setlist[itemize]{leftmargin=6mm,itemsep=2pt,topsep=3pt}
\setlist[enumerate]{leftmargin=7mm,itemsep=5pt,topsep=4pt}
\renewcommand{\arraystretch}{1.35}

\pagestyle{fancy}
\fancyhf{}
\fancyhead[L]{\small\textcolor{softgray}{Николь Саркисьянц · ЕГЭ}}
\fancyhead[R]{\small\textcolor{softgray}{Исследование функций}}
\fancyfoot[C]{\small\textcolor{softgray}{\thepage}}
\renewcommand{\headrulewidth}{0.3pt}
\renewcommand{\headrule}{%
  \hbox to\headwidth{%
    \color{accent}\leaders\hrule height \headrulewidth\hfill
  }%
}

\newcommand{\R}{\mathbb{R}}

\tikzset{
  startstop/.style={
    ellipse,
    draw=darkaccent,
    line width=0.9pt,
    minimum width=34mm,
    minimum height=10mm,
    align=center,
    text width=31mm
  },
  process/.style={
    rectangle,
    rounded corners=2mm,
    draw=accent,
    line width=0.9pt,
    minimum height=10mm,
    align=center,
    text width=92mm,
    inner sep=2.5mm,
    font=\small
  },
  decision/.style={
    diamond,
    aspect=2.7,
    draw=darkaccent,
    line width=0.9pt,
    align=center,
    text width=42mm,
    inner sep=1.2mm,
    font=\small
  },
  branch/.style={
    rectangle,
    rounded corners=2mm,
    draw=accent,
    line width=0.9pt,
    minimum height=11mm,
    align=center,
    text width=40mm,
    inner sep=2.2mm,
    font=\small
  },
  flowarrow/.style={
    -{Stealth[length=3mm,width=2mm]},
    line width=0.9pt,
    draw=darkaccent
  }
}

\begin{document}

% Основа содержания: доска занятия по исследованию функций.
% На ней зафиксирована схема:
% «Производная — Корни — Ось — Знаки — Ответ».
% :contentReference[oaicite:0]{index=0}

% ------------------------------------------------------------
% ТИТУЛЬНЫЙ ЛИСТ
% ------------------------------------------------------------

\thispagestyle{empty}

\begin{center}

\vspace*{30mm}

{\Large\textcolor{darkaccent}{\textbf{Чек-лист}}}\par

\vspace{7mm}

{\huge\textbf{Исследование функции с помощью производной}}\par

\vspace{5mm}

{\large
Монотонность · критические точки · экстремумы
}\par

\vspace{16mm}

{\large
Николь Саркисьянц · ЕГЭ
}\par

\vfill

{\large \today}\par

\vspace{8mm}

{\normalsize
Лёвин Артём Александрович эксклюзивно для Николь Саркисьянц
}\par

\vspace{18mm}

\end{center}

\begin{tikzpicture}[remember picture,overlay]
  \draw[
    accent,
    line width=1.1pt
  ]
    ([xshift=20mm,yshift=16mm]current page.south west)
    ..
    controls
    ([xshift=62mm,yshift=28mm]current page.south west)
    and
    ([xshift=-60mm,yshift=8mm]current page.south east)
    ..
    ([xshift=-20mm,yshift=18mm]current page.south east);
\end{tikzpicture}

\clearpage

% ------------------------------------------------------------
% ОСНОВНОЙ ЧЕК-ЛИСТ
% ------------------------------------------------------------

\section*{Основной чек-лист}

\subsection*{1. Что связывает производную и поведение функции}

Для внутренней точки области определения знак производной задаёт направление изменения функции:

\[
\begin{aligned}
f'(x)>0
&\quad\Longrightarrow\quad
f \text{ возрастает},
\\
f'(x)<0
&\quad\Longrightarrow\quad
f \text{ убывает}.
\end{aligned}
\]

Точка \(x_0\in D(f)\) называется \textbf{критической}, если
\[
f'(x_0)=0
\]
или производная в \(x_0\) не существует.

При исследовании экстремума принципиальна \textbf{смена знака} производной:

\[
\begin{array}{ccl}
+\ \longrightarrow\ -
&:&
\text{локальный максимум},
\\[1mm]
-\ \longrightarrow\ +
&:&
\text{локальный минимум},
\\[1mm]
+\ \longrightarrow\ +
\quad\text{или}\quad
-\ \longrightarrow\ -
&:&
\text{экстремума нет}.
\end{array}
\]

\textbf{Важно.}
Точки, в которых сама функция не определена, разделяют интервалы исследования, однако точками экстремума функции они быть не могут.

\subsection*{2. Памятка «ПКОЗА»}

\begin{center}

\begin{tabularx}{0.92\linewidth}{
  @{}
  >{\bfseries\color{darkaccent}}l
  X
  @{}
}
\toprule

П
&
Найдите производную \(f'(x)\).
\\

К
&
Найдите корни уравнения \(f'(x)=0\).
Отдельно учитывайте точки, где производная не существует, но сама функция определена.
\\

О
&
Нанесите критические точки и разрывы области определения на числовую ось.
\\

З
&
Определите знак \(f'(x)\) на каждом полученном интервале.
\\

А
&
Сформулируйте ответ:
интервалы возрастания и убывания,
точки максимума и минимума.
\\

\bottomrule
\end{tabularx}

\end{center}

\subsection*{3. Быстрый пример: квадратная функция}

Исследуем функцию

\[
f(x)=x^2-x-6.
\]

Её производная:

\[
f'(x)=2x-1.
\]

Находим критическую точку:

\[
2x-1=0
\Longrightarrow
x=\frac12.
\]

\begin{center}

\begin{tikzpicture}[x=2.5cm,y=1cm]

  \draw[-{Stealth[length=2.5mm]}]
    (-2,0)--(2.2,0)
    node[right] {$x$};

  \fill[darkaccent]
    (0,0) circle (1.8pt);

  \node[below=3pt]
    at (0,0)
    {$\frac12$};

  \node[
    above=4pt,
    text=darkaccent
  ]
    at (-1,0)
    {$f'(x)<0$};

  \node[
    above=4pt,
    text=darkaccent
  ]
    at (1,0)
    {$f'(x)>0$};

  \draw[
    accent,
    line width=0.9pt,
    -{Stealth[length=2mm]}
  ]
    (-1.4,-0.55)--(-0.15,-0.95);

  \draw[
    accent,
    line width=0.9pt,
    -{Stealth[length=2mm]}
  ]
    (0.15,-0.95)--(1.4,-0.55);

\end{tikzpicture}

\end{center}

Следовательно,

\[
f \text{ убывает на }
\left(-\infty;\frac12\right),
\]

\[
f \text{ возрастает на }
\left(\frac12;+\infty\right).
\]

В точке \(x=\frac12\) производная меняет знак с «\(-\)» на «\(+\)», поэтому там расположен локальный минимум.

\subsection*{4. Типовые ошибки}

\begin{itemize}

  \item
  забыть область определения до дифференцирования;

  \item
  принять точку разрыва функции за критическую точку;

  \item
  найти корни \(f'(x)\) и сразу назвать их экстремумами, не проверив смену знака;

  \item
  ошибиться со знаком производной на одном из интервалов;

  \item
  при произведении продифференцировать только один множитель;

  \item
  в дробно-рациональной функции потерять условие
  \[
  \text{знаменатель}\ne0.
  \]

\end{itemize}

\clearpage

% ------------------------------------------------------------
% РАЗОБРАННЫЕ КОНСТРУКЦИИ
% ------------------------------------------------------------

\section*{Разобранные конструкции с занятия}

\subsection*{1. Кубическая функция: два критических значения}

Рассмотрим

\[
f(x)
=
\frac13x^3
-
\frac12x^2
-
6x
+
5.
\]

Получаем:

\[
f'(x)
=
x^2-x-6
=
(x+2)(x-3).
\]

Критические значения:

\[
x=-2,
\qquad
x=3.
\]

Знак производной:

\[
\begin{array}{c|ccccc}
x
&
(-\infty,-2)
&
-2
&
(-2,3)
&
3
&
(3,+\infty)
\\
\midrule
f'(x)
&
+
&
0
&
-
&
0
&
+
\end{array}
\]

Поэтому функция возрастает на

\[
(-\infty,-2)
\quad\text{и}\quad
(3,+\infty),
\]

убывает на

\[
(-2,3).
\]

В точке \(x=-2\) расположен локальный максимум, в точке \(x=3\) --- локальный минимум.

\subsection*{2. Корень: сначала ОДЗ, затем производная}

Для функции

\[
f(x)
=
-x\sqrt{x}
+
3x
+
1
\]

область определения:

\[
D(f)=[0,+\infty).
\]

Удобно переписать

\[
x\sqrt{x}=x^{3/2}.
\]

Тогда при \(x>0\)

\[
(x\sqrt{x})'
=
\frac32\sqrt{x}.
\]

Следовательно,

\[
f'(x)
=
-\frac32\sqrt{x}
+
3.
\]

Ищем критическое значение:

\[
-\frac32\sqrt{x}+3=0,
\]

\[
\sqrt{x}=2,
\]

\[
x=4.
\]

На интервале

\[
(0,4)
\]

производная положительна, а на

\[
(4,+\infty)
\]

отрицательна.

Значит, \(x=4\) --- точка локального максимума.

\[
f(4)
=
-4\sqrt4+3\cdot4+1
=
-8+12+1
=
5.
\]

Точка максимума:

\[
(4,5).
\]

\textbf{Ключевая проверка.}
В произведении \(x\sqrt{x}\) нельзя учитывать производную только корня.

Удобный способ:

\[
x\sqrt{x}=x^{3/2}
\quad\Longrightarrow\quad
(x\sqrt{x})'
=
\frac32\sqrt{x}.
\]

\subsection*{3. Дробь: разрыв области определения обязан быть на оси}

Рассмотрим функцию

\[
f(x)
=
\frac{x^2+289}{x}.
\]

Сначала область определения:

\[
x\ne0.
\]

Упростим функцию:

\[
f(x)
=
x+\frac{289}{x}.
\]

Производная:

\[
f'(x)
=
1-\frac{289}{x^2}.
\]

Приведём к одной дроби:

\[
f'(x)
=
\frac{x^2-289}{x^2}.
\]

Корни производной:

\[
x^2-289=0,
\]

\[
x=\pm17.
\]

При этом точка

\[
x=0
\]

не входит в область определения.

Поэтому числовая ось разбивается на четыре интервала:

\[
(-\infty,-17),
\qquad
(-17,0),
\qquad
(0,17),
\qquad
(17,+\infty).
\]

Знаки производной соответственно:

\[
+,\quad -,\quad -,\quad +.
\]

Следовательно, локальный максимум достигается при

\[
x=-17.
\]

\[
f(-17)
=
-17-\frac{289}{17}
=
-34.
\]

Точка локального максимума:

\[
(-17,-34).
\]

Локальный минимум достигается при

\[
x=17.
\]

\[
f(17)
=
17+\frac{289}{17}
=
34.
\]

Точка локального минимума:

\[
(17,34).
\]

\subsection*{4. Произведение: полезно факторизовать производную}

Рассмотрим

\[
f(x)
=
(x+6)^2(x-10)+8.
\]

По правилу производной произведения:

\[
\begin{aligned}
f'(x)
&=
2(x+6)(x-10)
+
(x+6)^2
\\
&=
(x+6)
\bigl(
2x-20+x+6
\bigr)
\\
&=
(x+6)(3x-14).
\end{aligned}
\]

Факторизованный вид производной сразу даёт критические значения:

\[
x=-6,
\qquad
x=\frac{14}{3}.
\]

Такой вид значительно упрощает дальнейшее определение знаков производной.

\clearpage

% ------------------------------------------------------------
% БЛОК-СХЕМА
% ------------------------------------------------------------

\thispagestyle{fancy}

\begin{center}
{\Large
\textcolor{darkaccent}{
\textbf{Алгоритм: исследование функции по производной}
}
}
\end{center}

\vspace{2mm}

\begin{center}

\begin{tikzpicture}[
  node distance=4.5mm and 12mm
]

  \node[startstop]
    (start)
    {Начало};

  \node[
    process,
    below=of start
  ]
    (domain)
    {Найдите \(D(f)\) и отметьте точки разрыва области определения};

  \node[
    process,
    below=of domain
  ]
    (deriv)
    {Вычислите и упростите \(f'(x)\)};

  \node[
    process,
    below=of deriv
  ]
    (critical)
    {Найдите \(f'(x)=0\); учтите точки, где \(f'\) не существует, но \(f\) определена};

  \node[
    process,
    below=of critical
  ]
    (axis)
    {Отметьте на оси критические точки и разрывы; определите знаки \(f'(x)\)};

  \node[
    process,
    below=of axis
  ]
    (mono)
    {Запишите интервалы возрастания (\(f'>0\)) и убывания (\(f'<0\))};

  \node[
    decision,
    below=6mm of mono
  ]
    (change)
    {В критической точке знак \(f'\) меняется?};

  \node[
    branch,
    below=9mm of change,
    xshift=-47mm
  ]
    (ext)
    {Да:
    \(+\to-\) --- максимум;
    \(-\to+\) --- минимум};

  \node[
    branch,
    below=9mm of change,
    xshift=47mm
  ]
    (noext)
    {Нет:
    экстремума в этой критической точке нет};

  \node[
    process,
    below=31mm of change
  ]
    (answer)
    {Сформулируйте полный ответ: монотонность и точки экстремума};

  \node[
    startstop,
    below=of answer
  ]
    (finish)
    {Конец};

  \draw[flowarrow]
    (start)--(domain);

  \draw[flowarrow]
    (domain)--(deriv);

  \draw[flowarrow]
    (deriv)--(critical);

  \draw[flowarrow]
    (critical)--(axis);

  \draw[flowarrow]
    (axis)--(mono);

  \draw[flowarrow]
    (mono)--(change);

  \draw[flowarrow]
    (change.west)
    -|
    node[
      pos=0.20,
      above,
      font=\small
    ]
    {да}
    (ext.north);

  \draw[flowarrow]
    (change.east)
    -|
    node[
      pos=0.20,
      above,
      font=\small
    ]
    {нет}
    (noext.north);

  \draw[flowarrow]
    (ext.south)
    |-
    (answer.west);

  \draw[flowarrow]
    (noext.south)
    |-
    (answer.east);

  \draw[flowarrow]
    (answer)--(finish);

\end{tikzpicture}

\end{center}

\clearpage

% ------------------------------------------------------------
% ТРЕНИРОВОЧНЫЙ БЛОК
% ------------------------------------------------------------

\section*{Тренировочный блок: 5 заданий}

Для каждой функции:

\begin{enumerate}[label=\textbf{\arabic*.}]

  \item
  найдите область определения;

  \item
  исследуйте функцию на возрастание и убывание;

  \item
  укажите внутренние точки локального максимума и минимума, если они существуют.

\end{enumerate}

\vspace{3mm}

\begin{enumerate}[
  label=\textbf{Задание \arabic*.},
  leftmargin=5mm,
  labelindent=5mm,
  itemindent=5mm
]

  \item
  \[
  y=x^2-6x+5.
  \]

  \item
  \[
  y=-x^2+4x+1.
  \]

  \item
  \[
  y=
  \frac13x^3
  -
  x^2
  -
  3x
  +
  1.
  \]

  \item
  \[
  y=-x\sqrt{x}+6x+1.
  \]

  \item
  \[
  y=\frac{x^2+64}{x}.
  \]

\end{enumerate}

\vfill

\textcolor{softgray}{
\small
Проверка перед ответом:
ОДЗ
\(\rightarrow\)
производная
\(\rightarrow\)
корни
\(\rightarrow\)
ось
\(\rightarrow\)
знаки
\(\rightarrow\)
ответ.
}

\clearpage

% ------------------------------------------------------------
% ОТВЕТЫ
% ------------------------------------------------------------

\section*{Ответы к тренировочному блоку}

\small

\begin{tabularx}{\linewidth}{@{}c X@{}}

\toprule

\textbf{№}
&
\textbf{Ответ}
\\

\midrule

1
&
\(D=\R\).

Убывает на
\[
(-\infty,3),
\]
возрастает на
\[
(3,+\infty).
\]

Локальный минимум:
\[
(3,-4).
\]
\\

\addlinespace

2
&
\(D=\R\).

Возрастает на
\[
(-\infty,2),
\]
убывает на
\[
(2,+\infty).
\]

Локальный максимум:
\[
(2,5).
\]
\\

\addlinespace

3
&
\(D=\R\).

Возрастает на
\[
(-\infty,-1)\cup(3,+\infty),
\]
убывает на
\[
(-1,3).
\]

Локальный максимум:
\[
\left(-1,\frac83\right).
\]

Локальный минимум:
\[
(3,-8).
\]
\\

\addlinespace

4
&
\[
D=[0,+\infty).
\]

Возрастает на
\[
[0,16],
\]
убывает на
\[
(16,+\infty).
\]

Внутренний локальный максимум:
\[
(16,33).
\]
\\

\addlinespace

5
&
\[
D=\R\setminus\{0\}.
\]

Возрастает на
\[
(-\infty,-8)\cup(8,+\infty),
\]
убывает на
\[
(-8,0)\cup(0,8).
\]

Локальный максимум:
\[
(-8,-16).
\]

Локальный минимум:
\[
(8,16).
\]
\\

\bottomrule

\end{tabularx}

\normalsize

\vfill

\begin{center}
\textcolor{darkaccent}{
\textbf{
Главная идея:
знак производной описывает направление движения графика.
}
}
\end{center}

\end{document}